\documentclass[UTF8,fontset=none,11pt,a4paper]{ctexart}
\usepackage{ipara-handbook}
\usepackage{amsmath,booktabs,tabularx,hyperref}
\usetikzlibrary{arrows.meta,positioning,fit,calc}

\handbooksetup{OS-I02 fork + COW + Resource Reference}{408 · Operating System Integration}{Identity · Mapping · Reference · Divergence}

\begin{document}
\handbooktitle{\texttt{fork()} + COW + Resource Reference}{不要问“复制了什么”，要逐对象追踪复制、共享、引用与后续分化}{408 · OS-I02}

\begin{corebox}{Canonical Problem}
父进程调用 \code{fork()} 后，为什么父子既“像复制了一份进程”，又会暂时共享物理页和打开文件状态？随后一方执行写内存、\code{read/lseek/close} 或 \code{exec()} 时，哪些关系继续共享，哪些在什么事件上分化？
\[
\boxed{
\text{Parent State}
\xrightarrow{\texttt{fork}}
\text{New Child Identity}
+
\{\text{per-object relations}\}
\xrightarrow{\text{later events}}
\text{Selective Divergence}
}
\]
本册拒绝用“fork 就是把整个进程复制一遍”作为终点。真正稳定的做法是：\kw{对每一种对象分别判断 Copy / Share / Reference / Rebuild，以及以后哪一个事件会改变这条关系。}
\end{corebox}

\begin{methodbox}{Integration Ownership}
本册 Owns \kw{fork 后跨 Process / VM / File 的组合轨迹与分化时点}，不 Own 局部机制。
\begin{itemize}
  \item OS-01/02 Own process identity、fork/exec 生命周期、父子返回值与调度；
  \item OS-B02 + OS-04 Own address-space association、private mapping、COW fault/retry；
  \item OS-B03 + OS-06/07 Own fd binding、OFD/open instance、file object 与 namespace lifetime；
  \item 若 fork 后继续发生 I/O，再按具体操作调用 OS-05；若进入共享内存并发正确性，再调用 OS-03。
\end{itemize}
本 Canonical Process 的核心不依赖 I/O 或同步机制，因此它们不被强行塞进主线。
\end{methodbox}

\section{第一原则：fork 产生的是新的进程身份，不是“所有对象都复制”}

\code{fork()} 的核心语义是创建一个新的进程实例，使它从调用点附近继续执行，并继承父进程的一组环境关系。父子通过返回值区分当前控制流身份：
\[
\boxed{
\text{parent: child PID}
\qquad
\text{child: }0
}
\]

因此第一层必须独立出来：
\[
\boxed{\text{Process Identity}_{parent} \neq \text{Process Identity}_{child}}
\]
但“身份不同”并不推出所有底层资源立即拥有物理副本。fork 的设计恰好大量利用\kw{引用复制 + 延迟分化}。

\section{fork 后的关系向量：一次性把对象分层}

\begin{handbooklongtable}{L{2.7cm} Y Y Y}
\toprule
对象/关系 & fork 后稳定语义 & 常见高效实现 & 后续分化事件\\
\midrule
Process identity & child 是新的进程实例，PID/父子关系独立 & 新建内核管理记录 & exit/wait、独立调度生命周期\\
Execution context & 父子都从 fork 返回附近继续，但返回值等状态不同 & 初始化 child 的可恢复执行现场 & 两边之后各自执行不同指令流\\
Address space identity & 逻辑上是父子各自的地址空间环境 & 页表/region 关系复制或重建，私有页可暂时共享底层内容 & private write 触发 COW；exec 替换映像\\
Private writable page & 写入结果应彼此隔离 & 初始共享物理内容 + 写保护/等价 COW 元数据 & 第一次需要独立写时 mapping 分化\\
Explicit shared mapping & 共享语义继续存在 & 父子可继续引用同一共享 backing/page state & 不是 COW 隔离；由共享协议决定后续可见性\\
Descriptor table binding & child 得到父进程描述符绑定的语义副本 & 复制 fd binding 关系 & close/dup/exec close-on-exec 改变各自绑定\\
Open instance / OFD & 父子对应 fd 通常继续引用同一打开实例 & 增加引用关系，不复制 current offset & read/write/lseek 可改变共享 open-state\\
File object & 不因 fork 创建第二个文件对象 & 继续被相同 OFD 等引用 & namespace/open references 的变化由 FS lifecycle 决定\\
\bottomrule
\end{handbooklongtable}

这张表的核心不是记“复制/共享清单”，而是建立一个检查器：
\[
\boxed{\text{Identity} \to \text{Relation} \to \text{Owner} \to \text{Divergence Event}}
\]

\section{阶段一：fork 创建 child，并建立新的控制流身份}

父进程进入内核执行 fork 创建流程。OS-01/02 负责：
\begin{enumerate}
  \item 建立新的 child identity 和生命周期状态；
  \item 初始化 child 的执行现场，使其未来被调度后能从 fork 返回路径继续；
  \item 建立父子关系，并让父子在用户可观察返回值上不同；
  \item 把 child 送入合适的 Ready/Runnable 状态，真正何时执行由 scheduler 决定。
\end{enumerate}

因此：
\[
\boxed{\texttt{fork success} \not\Rightarrow \text{child immediately running}}
\]
创建完成只是让它成为可被调度的新进程。

\section{阶段二：地址空间是“逻辑独立 + 底层可暂时共享”}

OS-B02 把最容易混淆的两层拆开：
\begin{itemize}
  \item \kw{Address-space semantics}：父子私有内存写入之后应彼此隔离；
  \item \kw{Physical representation}：为了避免 fork 时复制所有页，多个 private mappings 可以暂时引用同一物理内容。
\end{itemize}

所以 COW 不是“fork 让父子共享内存”的最终语义，而是一种延迟兑现私有副本的 mechanism：
\[
\boxed{
\text{Private Semantics}
+
\text{Temporary Shared Representation}
\xrightarrow{\text{write}}
\text{Selective Physical Divergence}
}
\]

若题目出现 \code{MAP\_SHARED}/共享内存，则它本来就是 shared semantics，不能因为“fork 后两边都指到同一 frame”就把它误判成 COW private page。

\section{阶段三：文件关系不是“文件复制”，而是 binding 复制后继续引用同一 OFD}

对每个继承的 fd：
\[
\boxed{
\text{Parent fd binding}\to\text{OFD}
\qquad
\text{Child fd binding}\to\text{same OFD}
}
\]

这解释了一个经典结果：父子拥有\kw{各自的 fd 表绑定关系}，但对应绑定继续引用同一个打开实例，因此该 OFD 拥有的 current offset 等状态是共享的。

例如文件当前 offset 为 0，父进程先通过共享 OFD 成功读取 100 bytes，使 current offset 变为 100；如果 child 随后使用其继承 fd 做普通顺序 read，它从同一 OFD 的当前位置继续。真正判据不是“fd 数字是否相同”，而是：
\[
\boxed{\text{Do the bindings reference the same OFD?}}
\]

若父子后来各自重新 \code{open()} 同一路径，通常会得到不同 OFD；它们可以引用同一 File Object，却拥有独立 current offset。这正是 B03 的三层模型。

\section{四个后续事件：分化到底发生在哪一层}

\subsection{事件 A：父或子写 private page}

写保护/fault 进入 OS-04：
\[
\text{private write}
\to \text{COW decision}
\to \text{必要时准备私有页}
\to \text{更新写方 mapping}
\to \text{retry}
\]
只改变需要分化的 VM relation；PID、OFD、file object 不因此复制。

\subsection{事件 B：一方 \texttt{close(fd)}}

\code{close} 删除的是\kw{该进程的一条 descriptor binding}，并相应减少对打开实例的引用。另一进程自己的 binding 仍在时，它继续使用同一 OFD。于是：
\[
\boxed{\text{Child close} \not\Rightarrow \text{Parent fd invalid}}
\]

\subsection{事件 C：一方 \texttt{read()/lseek()}}

若父子 binding 指向同一 OFD，改变 current offset 的操作会影响另一方后续基于该 offset 的操作；但这不等于它们共享全部进程状态，也不等于文件对象被复制/修改。

\subsection{事件 D：一方 \texttt{exec()}}

\code{exec()} 的稳定边界是：
\[
\boxed{\text{same process identity, new program image}}
\]
进程身份不因 exec 自动变成另一个进程；原程序映像/地址空间内容被新程序替换。打开描述符是否保留要看 close-on-exec 等接口状态。于是 fork+exec 常见路径正好展示两种不同 operation：\kw{fork creates identity; exec replaces image}。

\section{贯穿母例：父子各做一次内存写和文件读}

假设父进程 P：
\begin{itemize}
  \item private page 中变量 \code{x=1}；
  \item fd=3 指向 OFD-A，current offset=0；
  \item OFD-A 指向文件对象 F。
\end{itemize}

P 调用 fork 得到 child C 后：
\begin{enumerate}
  \item P 与 C 的 process identity 已不同；
  \item \code{x} 所在 private page 的逻辑地址空间已属两边，但底层物理内容可以暂时共享；
  \item P:fd3 与 C:fd3 是两条 descriptor binding，二者继续引用 OFD-A；
  \item C 执行 \code{x=2}：只让 C 对该 private page 的 mapping 在 COW 时分化；P 仍看到自己的 private value；
  \item P 通过 fd3 读 100 bytes：OFD-A 的 current offset 按实际读取量推进；
  \item C 后续用继承 fd3 普通 read：从 OFD-A 的新 current offset 继续；COW 的内存分化并不会让 OFD 自动分化。
\end{enumerate}

这证明两个“共享/分化”轴彼此独立：
\[
\boxed{
\text{VM Mapping Relation}
\quad\perp\quad
\text{File Open-Reference Relation}
}
\]
一方发生 COW，不会顺带复制 OFD；一方 close fd，也不会改变另一方 private page。

\section{最容易错的五个绝对化说法}

\begin{handbooktable}{L{4.0cm} Y Y}
\toprule
错误压缩 & 为什么错 & 稳定判据\\
\midrule
“fork 把父进程完整复制一份” & 把语义复制、引用复制和物理复制混在一起 & 逐对象标 Copy / Share / Reference / Rebuild\\
“父子共享地址空间” & 私有映射语义要求后续写入隔离；只是底层页可暂时共享 & 问 mapping semantics 与 physical representation\\
“父子 fd 表共享” & fd binding 属于各自进程；共享的是绑定所引用的 OFD & 画 Parent fd $\to$ OFD 与 Child fd $\to$ OFD\\
“相同 fd 数字就共享 offset” & 两个独立 open 可用相同整数 fd 却指向不同 OFD & 只看是否同一 OFD\\
“exec 创建新进程” & exec 替换程序映像，不以创建新 PID 为定义 & 看 process identity 是否保持\\
\bottomrule
\end{handbooktable}

\section{扩展边界：哪些细节不要偷进 408 主模型}

\begin{itemize}
  \item 多线程进程调用 POSIX \code{fork()} 时 child 继承哪些线程，是接口级工程边界；只有题目明确涉及才展开，不反过来污染基础 process/VM/file 模型。
  \item COW 标志位放在哪里、页表层级怎样复制、是否原地复用唯一引用页、page size 多大，属于 OS-04 具体实现。
  \item inode/dentry/OFD 的精确内核引用计数与 journaling 回收时序属于 OS-06/07。
  \item fork 后真正何时父先跑还是子先跑由调度决定；没有题设不能根据代码顺序猜。
\end{itemize}

\section{独立验证：做 fork 综合题时逐列验收}

\begin{enumerate}
  \item \kw{Identity}：父子 PID/task identity 是否明确分开？
  \item \kw{Execution}：两条执行流分别从哪里继续，返回值如何区分？
  \item \kw{VM}：private/shared mapping 是否分开？物理共享是否只是实现表示？
  \item \kw{File}：fd binding、OFD、file object 是否三层分开？offset 属于哪一层？
  \item \kw{Event}：write/close/read/lseek/exec 分别改变哪条 relation，有没有错误“连带复制”？
  \item \kw{Scheduling}：child Ready 与 child Running 是否被分开？
\end{enumerate}

\begin{corebox}{Compression}
以后看到 fork，不再问一句“复制了什么”，而是画一张关系表：
\[
\boxed{
\text{New Identity}
+
\begin{cases}
\text{Execution Context: initialize/copy semantics}\\
\text{Private VM: logically separate, physically COW-share if allowed}\\
\text{Shared VM: keep shared semantics}\\
\text{fd binding: copy binding relation}\\
\text{OFD: referenced/shared}\\
\text{File Object: referenced, not duplicated}
\end{cases}
}
\]
然后对后续每个事件只修改它真正拥有的那一层。这个模型比“整进程复制”更接近 408 题目和真实 OS。
\end{corebox}

\end{document}
